\chapter{多模态模型}
\begin{quote}\small
\textit{本章摘要。}本章研究图像、文本、传感器和时间序列如何在同一任务中协同。先区分共同表示空间与真正的信息融合，再讨论对比学习、跨模态注意力、模态缺失和时间对齐，最后建立面向工业数据的评价方式。关键问题不是模态越多越好，而是每个模态是否提供了可对齐、可验证的互补信息。
\end{quote}
序列、视觉和语言模型各自提供证据后，还需要判断它们是否描述同一对象和同一阶段。本章先区分共享空间与信息融合，再处理对齐误差和模态缺失。联合诊断的类别标签与预警的未来事件标签必须分别定义；维修结论只有在形成后才能作为诊断证据，不能提前进入预警模型。

把照片和振动记录放在一起之前，先核对设备编号和采集时间。若照片来自维修后，振动来自维修前，直接配成一对就可能教错模型。共同嵌入空间用于寻找匹配记录，跨模态注意力让一种输入读取另一种输入，质量门控则在照片模糊或传感器失真时调整使用程度。这几种办法需要不同的训练数据。后文都用相机与振动联合诊断来说明：多出来的信息帮了什么忙，少一路输入时又该怎么办。

\section{不同模态的共同空间}

% auxiliary-resource-guide
本章末的“辅助学习材料”列出与核心方法对应的讲解和实践入口，可在阅读相关推导后，按所列小节对照学习。

\begin{figure}[htbp]
\centering
\begin{tikzpicture}[node distance=0.75cm and 0.9cm]
 \node[idi input, minimum width=1.8cm] (img) {图像\\$v_i$};
 \node[idi input, below=0.8cm of img, minimum width=1.8cm] (txt) {文本\\$t_i$};
 \node[idi process, right=of img, minimum width=1.8cm] (ev) {视觉编码器\\$f_v$};
 \node[idi process, right=of txt, minimum width=1.8cm] (et) {文本编码器\\$f_t$};
 \node[idi state, right=1.0cm of $(ev)!0.5!(et)$, minimum width=2.0cm] (space) {共同空间\\余弦相似度};
 \node[idi output, right=of space, minimum width=2.0cm] (loss) {对比损失\\匹配/非匹配};
 \draw[idi arrow] (img)--(ev); \draw[idi arrow] (txt)--(et);
 \draw[idi arrow] (ev)--(space); \draw[idi arrow] (et)--(space); \draw[idi arrow] (space)--(loss);
\end{tikzpicture}
\caption{双塔对比学习先把不同模态映射到共同空间，再用配对关系约束表示距离。}
\label{fig:multimodal-dual-encoder}
\end{figure}
图\ref{fig:multimodal-dual-encoder}的两路编码先独立计算，再在共同空间比较匹配关系；这种交互位置允许分别缓存表示，却未提供区域级联合推理。
模态（modality）是信息的表现形式，例如图像像素、文本词元和传感器数值；同一事件的不同模态仍可能有不同采样时间、噪声和缺失模式。编码器把一种原始输入变成特征，投影层再把不同宽度的特征映射到相同宽度；共同嵌入空间指这些输出可以用同一距离或内积比较，不要求原始输入格式相同。图文对比预训练（Contrastive Language--Image Pretraining，CLIP）通过图文对比学习获得可比较表示，BLIP-2 以轻量查询模块连接冻结的视觉与语言模型，Large Language and Vision Assistant (LLaVA) 通过视觉指令微调学习多模态交互\cite{radford2021clip,li2023blip2,liu2023llava}。BLIP-2 中的查询变换器称为 Querying Transformer (Q-Former)。三者分别强调对齐、表示连接和指令响应，并非同一损失的规模变化。
本章“对齐”指不同模态在对象、时间或表示上的对应，不同于上一章使回答符合指令与偏好的行为对齐；图文匹配良好不表示回答已符合规程。多模态模型首先要决定“对齐”发生在哪里。CLIP 采用图像编码器$f_v$和文本编码器$f_t$构成双塔结构，经过归一化后以温度参数$\tau$计算相似度：
\begin{equation}
 s_{ij}=\frac{f_v(v_i)^{\mathsf T}f_t(t_j)}{\tau}.
\end{equation}
批内对比目标让匹配图文相似度高于非匹配组合，因此 CLIP 能把“有裂纹”“轴承外圈剥落”等文本描述映射为分类原型。它的优点是图像和文本编码可以独立缓存，缺点是双塔交互浅，难以表达复杂问答和精确区域关系。

进入本章前，应能识别编码器输出的形状，理解内积、softmax和用标签计算损失；不必先训练一套视觉或语言大模型。先判断任务是寻找匹配记录、联合分类还是生成报告，再选双塔、融合头或语言连接模块。它们所需的监督分别可能是配对索引、类别标签和目标回答，不能只因输入都是图文就混用训练目标。

\paragraph{双塔的输入输出与一个两对样本的计算。}
设一批有$N=2$对图文，图像分支汇聚输出$2\times d$矩阵，文本分支也输出$2\times d$矩阵；行向量归一化是把每行除以其长度，要求非零向量。图像矩阵乘以文本矩阵的转置，再除以温度，得到$2\times2$得分表，其中对角线对应真实配对，非对角线是本批次用于比较的其他组合。设$d=2$，图文表示都依次为$(1,0),(0,1)$，取$\tau=1$，则得分表为$\left[\begin{smallmatrix}1&0\\0&1\end{smallmatrix}\right]$。每行正确配对概率为$e/(e+1)\approx0.731$，损失为$-\log0.731\approx0.313$。表示已经完全正交也不意味着softmax概率为一，因为它仍比较有限得分。

构造批次时，图文应按同一配对索引同步打乱；若单独重排文本，必须同步更新目标索引，不能继续把对角线当作真配对。训练时配对索引告诉模型哪个组合应更接近，梯度更新编码器或投影层；推理时不需要提供“正确配对”或重新计算训练损失。例如先缓存三种缺陷描述的向量，再把新图与它们比较，输出最相似的描述。所谓零样本分类只表示未用当前类别的标注训练专门分类头，不表示没有预训练；候选只有三类也不说明真实输入一定属于其中某类，需要另设拒识与校准。

BLIP-2 用冻结视觉编码器、冻结语言模型和轻量 Q-Former 连接二者。Q-Former 通过一组可学习查询从视觉特征中提取与文本任务相关的信息，再投影到语言模型输入空间；这样可以显著减少端到端训练成本，但视觉细节会受查询数量和投影瓶颈限制。LLaVA 类结构则通常把视觉编码器输出经过线性投影后拼接到语言模型上下文中，适合图像问答和报告生成。工业场景应把生成式描述与可定位模型分开验收：语言模型可以组织证据，不能凭文字流畅性证明缺陷确实位于某个像素区域。

\paragraph{连接语言模型时，向量还要经过什么步骤。}
假设视觉编码器产生$64\times128$个数，即64个区域、每区128维。Q-Former可用八个可学习查询读取这些区域，输出$8\times d_q$，再投影为$8\times d_{\rm lm}$，作为语言模型的条件输入。查询是训练得到的向量，不是八句人工问题。以LLaVA式直接投影为例，$128\times d_{\rm lm}$矩阵可把全部区域转成$64\times d_{\rm lm}$，与文字嵌入在序列维拼接。两条路线分别压缩视觉位置数与保留更多视觉位置，都会消耗语言上下文和计算预算。

训练图像描述时，语言部分仍预测目标回答的下一词元；推理时固定参数，给定图像和问题逐词生成。冻结视觉或语言权重不等于无需前向计算，也不等于从文字损失到连接模块之间无需传梯度。不同训练阶段可冻结不同部分，应报告实际设置。CLIP主要输出匹配分数，连接语言模型才获得文本生成路径；生成路径存在也不保证细裂纹的区域证据被保留。

对于传感器、图像和文本，常见做法是先分别得到$z_s,z_v,z_t$，再进行加权融合或跨模态注意力。若传感器序列的采样频率远高于图像帧率，可以用时间窗池化、事件 token 或 cross-attention 将高频状态压缩到视觉时间点；若某个模态临时失效，则训练时加入模态 dropout，并让质量门控$g_m$根据缺失率、噪声估计和传感器健康度动态调整贡献。多模态系统的关键不是把向量简单拼接，而是明确每种模态提供何种互补证据。

\begin{table}[htbp]
\centering
\caption{多模态模型的交互位置与适用任务}
\label{tab:multimodal-interaction}
\scriptsize
\setlength{\tabcolsep}{3pt}
\resizebox{\linewidth}{!}{%
\begin{tabular}{p{2.7cm}p{3.7cm}p{3.8cm}p{3.8cm}}
\hline
模型路线 & 交互机制 & 优势 & 工业风险与适用任务 \\
\hline
CLIP 双塔 & 图像--文本共享嵌入空间 & 编码可缓存、零样本分类方便 & 细粒度定位弱，适合筛查和检索 \\
BLIP-2 & Q-Former 从视觉特征提取查询 & 冻结大模型即可迁移 & 查询瓶颈可能丢失小缺陷，适合报告草拟 \\
LLaVA 类 & 视觉 token 投影到语言上下文 & 问答和多轮解释自然 & 生成内容需外部证据，不能代替定位器 \\
交叉注意力融合 & 查询模态读取键值模态 & 能表达局部、时间和条件关系 & 计算较重，适合联合诊断与根因分析 \\
\hline
\end{tabular}}
\end{table}

表\ref{tab:multimodal-interaction}对照共享嵌入、查询瓶颈和跨模态注意力，选型时须把编码成本与所需细粒度关系同时考虑。
以“相机图像加主轴振动诊断”为例，先由视觉编码器输出工件区域 token，由时序编码器输出若干转速同步的振动 token，再将维修记录编码为文本证据。模型可以采用两阶段决策：第一阶段用 CLIP 或轻量分类器筛选明显正常与明显异常样本，第二阶段仅对异常候选执行跨模态注意力和根因排序。训练标签应同时包含缺陷类别、严重度、发生时间和证据模态；如果只有最终维修结论，没有记录哪一帧图像或哪段频谱支持结论，模型即使获得较高准确率也难以审计。

联合损失可以写为
\begin{equation}
 \mathcal{L}=\lambda_{\rm cls}\mathcal{L}_{\rm cls}+\lambda_{\rm align}\mathcal{L}_{\rm InfoNCE}+\lambda_{\rm miss}\mathcal{L}_{\rm miss}+\lambda_{\rm rank}\mathcal{L}_{\rm rank},
\end{equation}
其中分类项约束业务输出，对比项约束共同空间，模态缺失项要求模型在删去某一路输入时仍保持可接受风险，排序项则约束根因候选的相对顺序。权重不能只按数值大小设置，应根据漏诊代价、人工复核容量和各模态的标签可信度进行敏感性分析。
图像、文本、音频和传感器序列具有不同的统计结构。多模态模型通常为每种模态建立编码器$z_m=f_m(x_m)$，再通过共享投影或跨模态注意力交互。对配对图文$(x_i,t_i)$，对比学习可最大化匹配样本相似度：
\begin{equation}
 \mathcal{L}_{\mathrm{NCE}}=-\frac{1}{N}\sum_i\log\frac{\exp(s(z_i^x,z_i^t)/\tau)}{\sum_j\exp(s(z_i^x,z_j^t)/\tau)}.
\end{equation}
这里$N$为配对样本数，$i,j=1,\ldots,N$，$z_i^x$与$z_i^v$均表示第$i$幅图像的单位范数表示，$s(a,b)=a^{\mathsf T}b$是不含温度的余弦相似度。因此$s_{ij}=s(z_i^v,z_j^t)/\tau$是已经缩放的对比得分，不能再次除以温度，且$\tau>0$。

\paragraph{融合方式。}
早期融合将原始特征拼接，晚期融合分别预测后再组合，中间融合在表示空间通过 cross-attention 交互。编码器--解码器架构还可以让一个模态条件化另一个模态的生成。工业场景中，时间同步、传感器缺失和模态质量差异往往比网络结构更决定系统效果。

\paragraph{拼接、加权和交叉读取的区别。}
若图像和振动分别输出三维、二维向量，拼接得到五维向量，可接$K\times5$分类矩阵预测$K$种状态；直接相加则不合法，必须先投影到同维度。晚期融合可以分别得到两个$K$类概率向量，再按权重组合，但要求类别含义和概率尺度可比较。中间交互则保留多个区域或时间片，让不同查询读取不同证据，不先压成各一个全局向量。

例如图片只显示表面颜色变化，振动给出周期冲击，两者联合可以帮助区分状态；若图片已丢掉裂纹细节，简单增加文字不能恢复真实像素。训练融合分类器还需要联合任务标签，只有成对图文并不自动提供故障类别、严重度或因果来源标签。前述联合损失是可选设计而非每个模型必备四项；缺失项和排序项应给出具体定义后才能复现实验。

\paragraph{多模态推理。}
图像问答、视觉定位和设备状态描述都可看作条件生成。应分别评估感知正确性、跨模态对齐和语言表达，避免只用语言流畅度掩盖视觉判断错误。

\paragraph{对比学习的温度参数。}
设图像和文本编码器输出归一化表示$z_i^v,z_i^t$，相似度取余弦相似度。温度$\tau$控制 softmax 的尖锐程度：
\begin{equation}
 s_{ij}=\frac{(z_i^v)^{\mathsf T}z_j^t}{\tau}.
\end{equation}
小温度会放大正负样本差异，梯度更集中于难负样本，但也可能使训练不稳定。对称的图到文、文到图损失通常比单向损失利用更多配对信息。

\paragraph{跨模态注意力。}
以文本为查询、视觉 patch 为键和值时，文本 token 可以选择性读取图像区域。为说明其计算含义，令$Q\in\mathbb R^{n_t\times d}$、$K\in\mathbb R^{n_v\times d}$、$V\in\mathbb R^{n_v\times d_v}$。矩阵$QK^{\mathsf T}$的第$(a,b)$项是第$a$个文本查询与第$b$个视觉键的内积；除以$\sqrt d$是为了在各维独立、方差为一时保持内积方差在可控范围。对每个文本位置$a$，定义
\[
\alpha_{ab}=\frac{\exp(q_a^{\mathsf T}k_b/\sqrt d)}{\sum_{r=1}^{n_v}\exp(q_a^{\mathsf T}k_r/\sqrt d)},\qquad \sum_b\alpha_{ab}=1.
\]
于是输出第$a$行为$h_a'=\sum_b\alpha_{ab}v_b$，它是视觉值向量的凸组合，而不是一个自动产生定位真值的坐标。若缺陷区域很小，$n_v$过小或下采样会使多个区域共享一个 token；若时间掩码错误，注意力仍会在错误窗口内给出看似合理的权重。实际部署应同时检查注意力之外的区域标注、遮挡实验和跨设备稳定性。
矩阵形式为
\begin{equation}
 H_{text}'=\softmax\left(\frac{Q_{text}K_{image}^{\mathsf T}}{\sqrt d}\right)V_{image}.
\end{equation}
此处$Q_{text}\in\R^{n_t\times d}$、$K_{image}\in\R^{n_v\times d}$、$V_{image}\in\R^{n_v\times d_v}$，$n_t,n_v$分别为文本与视觉 token 数。各行存放列向量的转置，softmax 沿视觉位置逐行计算，输出为$n_t\times d_v$矩阵。
这与简单拼接不同，因为每个文本位置可使用不同的视觉证据。多层交互可以支持图像描述、区域问答和视觉定位，但其对齐质量仍受配对数据噪声影响。

\paragraph{交叉注意力并没有把两个轴变成同一个轴。}
若两个文本查询读取三个视觉区域，则权重是$2\times3$矩阵，而不是$5\times5$。假设两行归一化权重分别是$(1/2,1/4,1/4)$和$(0.1,0.2,0.7)$，三个标量值为$2,4,10$，则输出为$4.5,8$；两个词元从同一张图读取了不同组合。输出行数跟查询数一致，要得到按视觉区域排列的结果，就需另设视觉查询路径或定位头，而不能直接把两个输出当成三个区域的标签。

\paragraph{模态缺失与质量加权。}
真实工业系统常出现相机被挡、传感器掉线或文本记录晚到。模型应显式接收“这一模态当前是否可用”的掩码$m_m$，并考虑
\begin{equation}
 z=\sum_m\frac{m_m\exp(a_m)}{\sum_r m_r\exp(a_r)}z_m.\label{eq:multimodal-quality-fusion}
\end{equation}
其中$m,r$遍历模态，$m_m\in\{0,1\}$表示第$m$个模态是否可用，$a_m$为有限的质量得分，各$z_m\in\R^{d_f}$须投影到相同维度，归一化权重是标量，融合结果$z\in\R^{d_f}$而非各模态维度之和。式\eqref{eq:multimodal-quality-fusion}要求$\sum_r m_r>0$；全部缺失时不计算这一比值，而应拒绝输出或调用预先验证的回退流程。训练时进行模态 dropout，能减少模型对某一模态的脆弱依赖，但必须保留真实缺失模式。

\paragraph{缺一路输入时，分母也必须改变。}
取两路质量得分$a_v=\log3,a_s=0$，完整输入时权重为$3/4,1/4$；若图像不可用，掩码变为$m_v=0,m_s=1$，权重应变为$0,1$，而不是仍把传感器乘以$1/4$。一个数值全零的图像可能是真实黑图，也可能是缺失填充值，必须用独立可用性掩码区分。模态dropout是训练时按设计暂时屏蔽某一路，教模型适应不完整输入，不是部署时随机扔掉可用观测。

共同空间的归一化同样需要边界处理：零向量不能除以自身长度，应标记为无效表示并检查输入或编码器；用小常数保护分母虽可防止除零，却不能让零向量成为单位方向。实现时应先清理缺失占位或跳过该分支，再参与融合；不能指望乘以零消除非数值（NaN），因为浮点运算中$0\times\mathrm{NaN}$仍为NaN。质量得分只用于相对加权，不自动等于该模态判断正确的校准概率。

一种可复现的缺失项是对至少保留一路的掩码模式重新前向计算，并对同一任务标签计算交叉熵，再在这些模式上平均。但若缺陷只能在图片中观察到，删图后标签在剩余输入下可能不可判定，不能强迫模型仍然自信答对；应允许不确定性或拒识，并按缺失模式分别验收。质量门控也不是裁判真相的机制，两路证据矛盾时还需检查数据来源。

\section{时间对齐、缺失鲁棒性与工业诊断}
\paragraph{时间对齐的多模态传感器。}
视觉帧、振动窗口和控制指令可能不在同一时间尺度。可以用时间编码、可学习延迟或 cross-attention 在邻近时间段内寻找对应信息。若把未来传感器值错误地对齐到当前图像，离线结果会明显偏高，部署时却失效。因此多模态数据集必须保存采集时间、传输延迟和有效时间区间。

设第$m$种模态在时间窗$W_i=[t_i^-,t_i^+]$内产生输入$x_i^{(m)}$，编码器输出$z_i^{(m)}
\in\mathbb R^d$。所谓配对，至少应满足设备标识一致、事件窗符合任务规定的重叠或滞后关系，并且作为输入的记录在预测时刻已经到达。监督标签可以在事后确认；训练时可用它计算损失，但不能把它或其派生字段提前作为模型输入。若任务要求同一事件的同步观测，而两个模态在校正允许滞后后仍无有效窗重叠，将它们强行作为正对会把时间错配当成监督噪声。对于存在未知延迟的系统，可在训练中显式搜索有限延迟集合$\Delta$，例如令
\[
\operatorname{sim}_i=\max_{\delta\in\Delta}\operatorname{sim}\bigl(z_i^{(v)},z_{i,\delta}^{(s)}\bigr),
\]
但最大化也可能选择偶然相似的错误事件，因此必须在留出设备和固定延迟范围上验证，而不能把对齐搜索结果直接解释为真实因果延迟。

\paragraph{按时钟核对一条样本。}
在10:00作出诊断时，振动窗口可以覆盖09:59:50至10:00，照片虽在09:59:58拍摄，若10:00:03才传到服务器，当前预测仍应把图像标记为不可用。10:30人工确认的缺陷类别可以作为这条样本的事后监督标签，却不是10:00已知的输入。把图像插值到最近传感器时点时，同样要检查插值是否用到了决策后的帧；按采集时间“对齐”不等于按可用时间因果合法。

将1000 Hz振动的十秒窗口与每秒一帧的图像配对时，振动原始输入有10000个采样点，不能仅因为都属于十秒就逐元素拼接。可先把振动分成十个因果一秒片段，每片编码为$d_s$维，再与对应图像区域交互；压缩也可能抹去短冲击，需比较不同片长。若模态间确有物理传播滞后，应保留有依据的滞后窗口，而非一律强求同一时刻。

\paragraph{图像与工艺曲线联合诊断。}
视觉编码器提取表面纹理，序列编码器提取温度和压力变化，再用融合层预测缺陷类型与严重度。消融实验至少包括仅视觉、仅时序、简单拼接和跨模态注意力四种设置。若融合模型只在完整数据上有效，还需报告相机缺失或曲线缺失时的退化曲线。

\paragraph{对齐、融合与协同的区别。}
对齐回答“不同模态是否指向同一个对象或事件”，融合回答“如何联合使用这些信息”，协同则进一步要求联合表示在任务上优于任何单一模态。三者不能混为一谈：图像和文本可以在共享空间中相近，却未必能完成精确定位；多个传感器可以简单拼接，却没有学到跨模态因果关系。

\begin{figure}[htbp]
\centering
\begin{tikzpicture}[>=Latex, node distance=0.8cm, every node/.style={font=\small}]
 \node[draw, rounded corners, fill=blue!10, minimum width=2.4cm] (vision) {视觉编码器};
 \node[draw, rounded corners, fill=orange!12, right=2cm of vision, minimum width=2.4cm] (text) {文本编码器};
 \node[draw, rounded corners, fill=green!12, below=1.1cm of vision, minimum width=2.4cm] (sensor) {传感器编码器};
 \node[draw, rounded corners, fill=yellow!20, right=2cm of sensor, minimum width=2.4cm] (fusion) {对齐/融合层};
 \node[draw, rounded corners, fill=red!10, right=2cm of fusion, minimum width=2.4cm] (out) {诊断或生成};
 \draw[->, thick] (vision)--(fusion); \draw[->, thick] (text)--(fusion); \draw[->, thick] (sensor)--(fusion); \draw[->, thick] (fusion)--(out);
\end{tikzpicture}
\caption{工业多模态系统中，编码、对齐、融合和任务输出是不同层次。}
\label{fig:multimodal-industrial-fusion}
\end{figure}
图\ref{fig:multimodal-industrial-fusion}将各模态编码器与融合层分开，说明原始输入先形成各自表示，再由对齐与融合环节为诊断或生成提供联合证据。

\paragraph{对称对比损失。}
图到文和文到图两个方向都应参与训练。令$S_{ij}=s_{ij}$为已经除以温度的得分矩阵，且$i,j=1,\ldots,N$，则对称损失为
\begin{equation}
 \mathcal{L}=\frac12\left[-\frac1N\sum_i\log\frac{e^{S_{ii}}}{\sum_j e^{S_{ij}}}
 -\frac1N\sum_i\log\frac{e^{S_{ii}}}{\sum_j e^{S_{ji}}}\right].
\end{equation}
第一项固定图像$i$、沿第$i$行在文本间归一化；第二项固定文本$i$、沿第$i$列在图像间归一化，两者都把对角位置作为目标。因此第二项分母使用$S_{ji}$而非$S_{ij}$。这里正样本指匹配的图文对，负样本指作为不匹配候选的图文对，与设备“正常／故障”的类别无关。若$N=1$且没有额外负样本，两方向的分子分母相同，损失恒为零，也没有对比分数梯度；因此损失小不一定表示已学到匹配。批次中的负样本并不都是真负样本：两个不同时间的图像可能描述同一设备状态，两个文本也可能共享相同故障原因。因此对比学习需要检查假负例（false negative，即实际匹配却被当作负例的组合），并在必要时使用组标签、软标签或难负样本过滤。

\paragraph{模态不平衡与缺失鲁棒性。}
某一模态信噪比更高时，模型可能忽略其他模态；训练损失下降并不表示融合真正发生。可以通过模态 dropout、分支独立监督、梯度平衡和质量门控减轻这一问题。测试时应绘制从完整输入到逐模态缺失的性能曲线，并区分随机缺失、连续掉线、延迟到达和系统性遮挡。

\paragraph{时间对齐与事件粒度。}
图像帧、振动窗口和控制指令的时间粒度不同。对齐时要明确事件发生时间、采集时间、传输时间和标签确认时间。窗口重叠会增加样本条目数，但不等于增加同样多的独立信息，还可能把同一故障事件复制到训练和测试两侧。应以事件或设备为单位划分，并在报告中说明窗口长度、步长和允许的预测提前量。

\paragraph{一个能区分互补信息与模型规模的评价。}
先固定同一批测试事件、同一分类标签与决策时间，在相同调参预算下比较仅图像、仅振动、拼接与交叉注意力，再按照片模糊、传感器掉线和两路冲突分别统计。假设一百个事件中图像答对八十个、振动答对七十五个，这不足以推出融合能答对九十个：两路错误可能高度重合，也可能互补，须逐事件比较预测。还可打乱图像与振动的配对关系作对照；若性能不变，可能没有利用对应关系，但也可能任务本来不需要它，不能仅据此断定实现错误。

用于报告生成时，至少把“识别对了什么”“证据对应哪帧、哪段信号”“文字是否忠于证据”分别计分。正确类别加上编造的位置不算定位正确；只会在完整输入下生成流畅报告，也不能代替缺失模式下的拒识验证。所有示意数值用于解释评价方法，不代表本书已完成这些模型的实测比较。

\paragraph{多模态评价。}
多模态系统应分别报告单模态基线、简单拼接、注意力融合和缺失模态结果。分类任务可检查跨模态遮挡后的鲁棒性；定位任务需要区域级准确率；生成任务需要事实一致性和证据覆盖率。最重要的消融不是把某一层删掉，而是回答“哪一种模态信息在什么工况下真正改变了决策”。

多模态学习的共同嵌入、对比目标与跨模态注意力可用下列材料互补复习，分别侧重表示学习和工程实现。
\section*{辅助学习材料}
\begingroup
\emergencystretch=3em
\begin{itemize}
\item \textbf{Contrastive Representation Learning}；作者/平台：Lilian Weng；英文；技术教程。重点看 contrastive loss、InfoNCE、temperature 和 negative sampling，帮助理解本章的 CLIP 双塔、对齐空间与模态缺失风险。\href{https://lilianweng.github.io/posts/2021-05-31-contrastive/}{在线阅读}。
\item \textbf{Natural language image search with a Dual Encoder}；作者/平台：Khalid Salama / Keras；英文；代码教程。阅读双编码器、投影层、对比目标和文本检索图像部分，理解共享表示空间。示例使用 Keras 2，运行前需核对依赖版本。\href{https://keras.io/examples/vision/nl\_image\_search/}{在线阅读}。
\item \textbf{多模态学习综述 (MultiModal Learning)}；知乎；中文技术解说。区分表示、对齐与融合，核对不同模态的数据配对方式；标题与主题据必应索引，原页受限。\href{https://zhuanlan.zhihu.com/p/582878508}{在线阅读}。
\item \textbf{一文彻底搞懂多模态：模态表示、多模态融合、跨模态对齐（索引标题节选）}；CSDN；中文博客。比较早期与后期融合，另行检查缺失模态下的评价；标题与主题据必应索引，原页受限，完整标题未核实。\href{https://blog.csdn.net/star\_nwe/article/details/143416379}{在线阅读}。
\item \textbf{多模态 之 CLIP ｜ 1、CLIP 网络结构精讲}；上传者：Enzo\_Mi；bilibili；中文技术讲解。对照本章图文对齐，画出图像编码器、文本编码器和相似度计算，比较对比学习与融合；不要把图文匹配直接等同于生成能力。2026年10月4日公开元数据播放4,447次；优先选择结构讲解而非成品展示。视频落地页显示验证码，未绕过；标题、上传者和计数据官方公开元数据核对，未逐帧观看。\href{https://www.bilibili.com/video/BV1BeZFBGEhe/}{观看视频}。
\end{itemize}
\endgroup
\paragraph{自测与参考答案。}
沿用10:00诊断例：09:59:58拍摄、10:00:03才到达的照片，与10:30确认的缺陷类别，哪一个可用于10:00预测，哪一个可用于事后训练监督？参考答案：照片尚未到达，不能作为该时刻输入，应标记缺失；确认类别可作为监督标签，但不能回填成当时已知的特征。若图像缺失而传感器可用，质量权重须重新归一化为$0,1$；若两路均缺失，分母为零，应进入拒识或已验证的回退流程。

\paragraph{本章小结。}
图像、振动和文字的采集频率不同，质量也不同，还可能缺失。联合使用时，要说明每一路怎样编码、按什么时间和设备配对、相互矛盾时听谁的，以及少一路输入时还能完成什么任务。最后在独立数据上分别测试完整输入和缺失输入，才能判断复杂的融合结构是否真的有帮助。
